% 原创TikZ重绘；层次布局参考Bishop与Bishop（2024）第12章。
\begin{tikzpicture}[x=1mm,y=1mm,font=\figurefont\footnotesize,
 every node/.style={text=NNInk},
 stage/.style={nn operator,minimum height=9mm,inner sep=2pt},
 norm/.style={stage,minimum height=7mm},
 wire/.style={nn flow,draw=NNWire,rounded corners=1.5mm},
 aux/.style={draw=NNLine,line width=.85pt,densely dashed},
 sum/.style={nn pointwise,minimum size=5mm}]

\node[] (r1tok0) at (20,0) {猫};
\node[inner sep=0pt,minimum width=12mm,minimum height=3mm] (r1e0) at (20,12) {};
\filldraw[draw=NNInput,fill=NNInput!18,line width=.5pt] (14.0,10.5) rectangle ++(3.0,3);
\filldraw[draw=NNInput,fill=NNInput!18,line width=.5pt] (17.0,10.5) rectangle ++(3.0,3);
\filldraw[draw=NNInput,fill=NNInput!18,line width=.5pt] (20.0,10.5) rectangle ++(3.0,3);
\filldraw[draw=NNInput,fill=NNInput!18,line width=.5pt] (23.0,10.5) rectangle ++(3.0,3);
\draw[wire] (r1tok0.north)--(r1e0.south);
\node[] (r1plus0) at (20,18) {$+$};
\node[inner sep=0pt,minimum width=12mm,minimum height=3mm] (r1pos0) at (20,24) {};
\filldraw[draw=NNInput,fill=NNInput!18,line width=.5pt] (14.0,22.5) rectangle ++(3.0,3);
\filldraw[draw=NNInput,fill=NNInput!18,line width=.5pt] (17.0,22.5) rectangle ++(3.0,3);
\filldraw[draw=NNInput,fill=NNInput!18,line width=.5pt] (20.0,22.5) rectangle ++(3.0,3);
\filldraw[draw=NNInput,fill=NNInput!18,line width=.5pt] (23.0,22.5) rectangle ++(3.0,3);
\draw[aux] (13,8) rectangle (27,28);
\node[inner sep=0pt,minimum width=12mm,minimum height=3mm] (r1h00) at (20,39) {};
\filldraw[draw=NNInput,fill=NNInput!18,line width=.5pt] (14.0,37.5) rectangle ++(3.0,3);
\filldraw[draw=NNInput,fill=NNInput!18,line width=.5pt] (17.0,37.5) rectangle ++(3.0,3);
\filldraw[draw=NNInput,fill=NNInput!18,line width=.5pt] (20.0,37.5) rectangle ++(3.0,3);
\filldraw[draw=NNInput,fill=NNInput!18,line width=.5pt] (23.0,37.5) rectangle ++(3.0,3);
\draw[wire] (20,28)--(r1h00.south);
\node[] (r1tok1) at (37,0) {[MASK]};
\node[inner sep=0pt,minimum width=12mm,minimum height=3mm] (r1e1) at (37,12) {};
\filldraw[draw=NNInput,fill=NNInput!18,line width=.5pt] (31.0,10.5) rectangle ++(3.0,3);
\filldraw[draw=NNInput,fill=NNInput!18,line width=.5pt] (34.0,10.5) rectangle ++(3.0,3);
\filldraw[draw=NNInput,fill=NNInput!18,line width=.5pt] (37.0,10.5) rectangle ++(3.0,3);
\filldraw[draw=NNInput,fill=NNInput!18,line width=.5pt] (40.0,10.5) rectangle ++(3.0,3);
\draw[wire] (r1tok1.north)--(r1e1.south);
\node[] (r1plus1) at (37,18) {$+$};
\node[inner sep=0pt,minimum width=12mm,minimum height=3mm] (r1pos1) at (37,24) {};
\filldraw[draw=NNInput,fill=NNInput!18,line width=.5pt] (31.0,22.5) rectangle ++(3.0,3);
\filldraw[draw=NNInput,fill=NNInput!18,line width=.5pt] (34.0,22.5) rectangle ++(3.0,3);
\filldraw[draw=NNInput,fill=NNInput!18,line width=.5pt] (37.0,22.5) rectangle ++(3.0,3);
\filldraw[draw=NNInput,fill=NNInput!18,line width=.5pt] (40.0,22.5) rectangle ++(3.0,3);
\draw[aux] (30,8) rectangle (44,28);
\node[inner sep=0pt,minimum width=12mm,minimum height=3mm] (r1h01) at (37,39) {};
\filldraw[draw=NNInput,fill=NNInput!18,line width=.5pt] (31.0,37.5) rectangle ++(3.0,3);
\filldraw[draw=NNInput,fill=NNInput!18,line width=.5pt] (34.0,37.5) rectangle ++(3.0,3);
\filldraw[draw=NNInput,fill=NNInput!18,line width=.5pt] (37.0,37.5) rectangle ++(3.0,3);
\filldraw[draw=NNInput,fill=NNInput!18,line width=.5pt] (40.0,37.5) rectangle ++(3.0,3);
\draw[wire] (37,28)--(r1h01.south);
\node[] (r1tok2) at (54,0) {狗};
\node[inner sep=0pt,minimum width=12mm,minimum height=3mm] (r1e2) at (54,12) {};
\filldraw[draw=NNInput,fill=NNInput!18,line width=.5pt] (48.0,10.5) rectangle ++(3.0,3);
\filldraw[draw=NNInput,fill=NNInput!18,line width=.5pt] (51.0,10.5) rectangle ++(3.0,3);
\filldraw[draw=NNInput,fill=NNInput!18,line width=.5pt] (54.0,10.5) rectangle ++(3.0,3);
\filldraw[draw=NNInput,fill=NNInput!18,line width=.5pt] (57.0,10.5) rectangle ++(3.0,3);
\draw[wire] (r1tok2.north)--(r1e2.south);
\node[] (r1plus2) at (54,18) {$+$};
\node[inner sep=0pt,minimum width=12mm,minimum height=3mm] (r1pos2) at (54,24) {};
\filldraw[draw=NNInput,fill=NNInput!18,line width=.5pt] (48.0,22.5) rectangle ++(3.0,3);
\filldraw[draw=NNInput,fill=NNInput!18,line width=.5pt] (51.0,22.5) rectangle ++(3.0,3);
\filldraw[draw=NNInput,fill=NNInput!18,line width=.5pt] (54.0,22.5) rectangle ++(3.0,3);
\filldraw[draw=NNInput,fill=NNInput!18,line width=.5pt] (57.0,22.5) rectangle ++(3.0,3);
\draw[aux] (47,8) rectangle (61,28);
\node[inner sep=0pt,minimum width=12mm,minimum height=3mm] (r1h02) at (54,39) {};
\filldraw[draw=NNInput,fill=NNInput!18,line width=.5pt] (48.0,37.5) rectangle ++(3.0,3);
\filldraw[draw=NNInput,fill=NNInput!18,line width=.5pt] (51.0,37.5) rectangle ++(3.0,3);
\filldraw[draw=NNInput,fill=NNInput!18,line width=.5pt] (54.0,37.5) rectangle ++(3.0,3);
\filldraw[draw=NNInput,fill=NNInput!18,line width=.5pt] (57.0,37.5) rectangle ++(3.0,3);
\draw[wire] (54,28)--(r1h02.south);
\node[stage,minimum width=51mm,minimum height=10mm] (r1layer1) at (37.0,57) {第1层：双向读取};
\node[stage,minimum width=51mm,minimum height=10mm] (r1layerL) at (37.0,85) {第$L$层：双向读取};
\draw[wire] (r1h00.north)--(20,52);
\draw[wire] (20,62)--(20,67);
\node[] (r1ellipsis0) at (20,71) {$\vdots$};
\draw[wire] (20,75)--(20,80);
\node[inner sep=0pt,minimum width=12mm,minimum height=3mm] (r1hL0) at (20,102) {};
\filldraw[draw=NNOutput,fill=NNOutput!18,line width=.5pt] (14.0,100.5) rectangle ++(3.0,3);
\filldraw[draw=NNOutput,fill=NNOutput!18,line width=.5pt] (17.0,100.5) rectangle ++(3.0,3);
\filldraw[draw=NNOutput,fill=NNOutput!18,line width=.5pt] (20.0,100.5) rectangle ++(3.0,3);
\filldraw[draw=NNOutput,fill=NNOutput!18,line width=.5pt] (23.0,100.5) rectangle ++(3.0,3);
\draw[wire] (20,90)--(r1hL0.south);
\draw[wire] (r1h01.north)--(37,52);
\draw[wire] (37,62)--(37,67);
\node[] (r1ellipsis1) at (37,71) {$\vdots$};
\draw[wire] (37,75)--(37,80);
\node[inner sep=0pt,minimum width=12mm,minimum height=3mm] (r1hL1) at (37,102) {};
\filldraw[draw=NNOutput,fill=NNOutput!18,line width=.5pt] (31.0,100.5) rectangle ++(3.0,3);
\filldraw[draw=NNOutput,fill=NNOutput!18,line width=.5pt] (34.0,100.5) rectangle ++(3.0,3);
\filldraw[draw=NNOutput,fill=NNOutput!18,line width=.5pt] (37.0,100.5) rectangle ++(3.0,3);
\filldraw[draw=NNOutput,fill=NNOutput!18,line width=.5pt] (40.0,100.5) rectangle ++(3.0,3);
\draw[wire] (37,90)--(r1hL1.south);
\draw[wire] (r1h02.north)--(54,52);
\draw[wire] (54,62)--(54,67);
\node[] (r1ellipsis2) at (54,71) {$\vdots$};
\draw[wire] (54,75)--(54,80);
\node[inner sep=0pt,minimum width=12mm,minimum height=3mm] (r1hL2) at (54,102) {};
\filldraw[draw=NNOutput,fill=NNOutput!18,line width=.5pt] (48.0,100.5) rectangle ++(3.0,3);
\filldraw[draw=NNOutput,fill=NNOutput!18,line width=.5pt] (51.0,100.5) rectangle ++(3.0,3);
\filldraw[draw=NNOutput,fill=NNOutput!18,line width=.5pt] (54.0,100.5) rectangle ++(3.0,3);
\filldraw[draw=NNOutput,fill=NNOutput!18,line width=.5pt] (57.0,100.5) rectangle ++(3.0,3);
\draw[wire] (54,90)--(r1hL2.south);
\node[stage,draw=NNOutput,fill=NNOutput!18,minimum width=22mm,minimum height=8mm] (r1head) at (37,117) {MLM头};
\draw[wire] (r1hL1.north)--(r1head.south);
\node[] (r1target) at (37,130) {追};
\draw[wire] (r1head.north)--(r1target.south);
\node[nn heading] (r1sample) at (37,-12) {遮蔽样本1};
\node[] (r2tok0) at (77,0) {[MASK]};
\node[inner sep=0pt,minimum width=12mm,minimum height=3mm] (r2e0) at (77,12) {};
\filldraw[draw=NNInput,fill=NNInput!18,line width=.5pt] (71.0,10.5) rectangle ++(3.0,3);
\filldraw[draw=NNInput,fill=NNInput!18,line width=.5pt] (74.0,10.5) rectangle ++(3.0,3);
\filldraw[draw=NNInput,fill=NNInput!18,line width=.5pt] (77.0,10.5) rectangle ++(3.0,3);
\filldraw[draw=NNInput,fill=NNInput!18,line width=.5pt] (80.0,10.5) rectangle ++(3.0,3);
\draw[wire] (r2tok0.north)--(r2e0.south);
\node[] (r2plus0) at (77,18) {$+$};
\node[inner sep=0pt,minimum width=12mm,minimum height=3mm] (r2pos0) at (77,24) {};
\filldraw[draw=NNInput,fill=NNInput!18,line width=.5pt] (71.0,22.5) rectangle ++(3.0,3);
\filldraw[draw=NNInput,fill=NNInput!18,line width=.5pt] (74.0,22.5) rectangle ++(3.0,3);
\filldraw[draw=NNInput,fill=NNInput!18,line width=.5pt] (77.0,22.5) rectangle ++(3.0,3);
\filldraw[draw=NNInput,fill=NNInput!18,line width=.5pt] (80.0,22.5) rectangle ++(3.0,3);
\draw[aux] (70,8) rectangle (84,28);
\node[inner sep=0pt,minimum width=12mm,minimum height=3mm] (r2h00) at (77,39) {};
\filldraw[draw=NNInput,fill=NNInput!18,line width=.5pt] (71.0,37.5) rectangle ++(3.0,3);
\filldraw[draw=NNInput,fill=NNInput!18,line width=.5pt] (74.0,37.5) rectangle ++(3.0,3);
\filldraw[draw=NNInput,fill=NNInput!18,line width=.5pt] (77.0,37.5) rectangle ++(3.0,3);
\filldraw[draw=NNInput,fill=NNInput!18,line width=.5pt] (80.0,37.5) rectangle ++(3.0,3);
\draw[wire] (77,28)--(r2h00.south);
\node[] (r2tok1) at (94,0) {追};
\node[inner sep=0pt,minimum width=12mm,minimum height=3mm] (r2e1) at (94,12) {};
\filldraw[draw=NNInput,fill=NNInput!18,line width=.5pt] (88.0,10.5) rectangle ++(3.0,3);
\filldraw[draw=NNInput,fill=NNInput!18,line width=.5pt] (91.0,10.5) rectangle ++(3.0,3);
\filldraw[draw=NNInput,fill=NNInput!18,line width=.5pt] (94.0,10.5) rectangle ++(3.0,3);
\filldraw[draw=NNInput,fill=NNInput!18,line width=.5pt] (97.0,10.5) rectangle ++(3.0,3);
\draw[wire] (r2tok1.north)--(r2e1.south);
\node[] (r2plus1) at (94,18) {$+$};
\node[inner sep=0pt,minimum width=12mm,minimum height=3mm] (r2pos1) at (94,24) {};
\filldraw[draw=NNInput,fill=NNInput!18,line width=.5pt] (88.0,22.5) rectangle ++(3.0,3);
\filldraw[draw=NNInput,fill=NNInput!18,line width=.5pt] (91.0,22.5) rectangle ++(3.0,3);
\filldraw[draw=NNInput,fill=NNInput!18,line width=.5pt] (94.0,22.5) rectangle ++(3.0,3);
\filldraw[draw=NNInput,fill=NNInput!18,line width=.5pt] (97.0,22.5) rectangle ++(3.0,3);
\draw[aux] (87,8) rectangle (101,28);
\node[inner sep=0pt,minimum width=12mm,minimum height=3mm] (r2h01) at (94,39) {};
\filldraw[draw=NNInput,fill=NNInput!18,line width=.5pt] (88.0,37.5) rectangle ++(3.0,3);
\filldraw[draw=NNInput,fill=NNInput!18,line width=.5pt] (91.0,37.5) rectangle ++(3.0,3);
\filldraw[draw=NNInput,fill=NNInput!18,line width=.5pt] (94.0,37.5) rectangle ++(3.0,3);
\filldraw[draw=NNInput,fill=NNInput!18,line width=.5pt] (97.0,37.5) rectangle ++(3.0,3);
\draw[wire] (94,28)--(r2h01.south);
\node[] (r2tok2) at (111,0) {狗};
\node[inner sep=0pt,minimum width=12mm,minimum height=3mm] (r2e2) at (111,12) {};
\filldraw[draw=NNInput,fill=NNInput!18,line width=.5pt] (105.0,10.5) rectangle ++(3.0,3);
\filldraw[draw=NNInput,fill=NNInput!18,line width=.5pt] (108.0,10.5) rectangle ++(3.0,3);
\filldraw[draw=NNInput,fill=NNInput!18,line width=.5pt] (111.0,10.5) rectangle ++(3.0,3);
\filldraw[draw=NNInput,fill=NNInput!18,line width=.5pt] (114.0,10.5) rectangle ++(3.0,3);
\draw[wire] (r2tok2.north)--(r2e2.south);
\node[] (r2plus2) at (111,18) {$+$};
\node[inner sep=0pt,minimum width=12mm,minimum height=3mm] (r2pos2) at (111,24) {};
\filldraw[draw=NNInput,fill=NNInput!18,line width=.5pt] (105.0,22.5) rectangle ++(3.0,3);
\filldraw[draw=NNInput,fill=NNInput!18,line width=.5pt] (108.0,22.5) rectangle ++(3.0,3);
\filldraw[draw=NNInput,fill=NNInput!18,line width=.5pt] (111.0,22.5) rectangle ++(3.0,3);
\filldraw[draw=NNInput,fill=NNInput!18,line width=.5pt] (114.0,22.5) rectangle ++(3.0,3);
\draw[aux] (104,8) rectangle (118,28);
\node[inner sep=0pt,minimum width=12mm,minimum height=3mm] (r2h02) at (111,39) {};
\filldraw[draw=NNInput,fill=NNInput!18,line width=.5pt] (105.0,37.5) rectangle ++(3.0,3);
\filldraw[draw=NNInput,fill=NNInput!18,line width=.5pt] (108.0,37.5) rectangle ++(3.0,3);
\filldraw[draw=NNInput,fill=NNInput!18,line width=.5pt] (111.0,37.5) rectangle ++(3.0,3);
\filldraw[draw=NNInput,fill=NNInput!18,line width=.5pt] (114.0,37.5) rectangle ++(3.0,3);
\draw[wire] (111,28)--(r2h02.south);
\node[stage,minimum width=51mm,minimum height=10mm] (r2layer1) at (94.0,57) {第1层：双向读取};
\node[stage,minimum width=51mm,minimum height=10mm] (r2layerL) at (94.0,85) {第$L$层：双向读取};
\draw[wire] (r2h00.north)--(77,52);
\draw[wire] (77,62)--(77,67);
\node[] (r2ellipsis0) at (77,71) {$\vdots$};
\draw[wire] (77,75)--(77,80);
\node[inner sep=0pt,minimum width=12mm,minimum height=3mm] (r2hL0) at (77,102) {};
\filldraw[draw=NNOutput,fill=NNOutput!18,line width=.5pt] (71.0,100.5) rectangle ++(3.0,3);
\filldraw[draw=NNOutput,fill=NNOutput!18,line width=.5pt] (74.0,100.5) rectangle ++(3.0,3);
\filldraw[draw=NNOutput,fill=NNOutput!18,line width=.5pt] (77.0,100.5) rectangle ++(3.0,3);
\filldraw[draw=NNOutput,fill=NNOutput!18,line width=.5pt] (80.0,100.5) rectangle ++(3.0,3);
\draw[wire] (77,90)--(r2hL0.south);
\draw[wire] (r2h01.north)--(94,52);
\draw[wire] (94,62)--(94,67);
\node[] (r2ellipsis1) at (94,71) {$\vdots$};
\draw[wire] (94,75)--(94,80);
\node[inner sep=0pt,minimum width=12mm,minimum height=3mm] (r2hL1) at (94,102) {};
\filldraw[draw=NNOutput,fill=NNOutput!18,line width=.5pt] (88.0,100.5) rectangle ++(3.0,3);
\filldraw[draw=NNOutput,fill=NNOutput!18,line width=.5pt] (91.0,100.5) rectangle ++(3.0,3);
\filldraw[draw=NNOutput,fill=NNOutput!18,line width=.5pt] (94.0,100.5) rectangle ++(3.0,3);
\filldraw[draw=NNOutput,fill=NNOutput!18,line width=.5pt] (97.0,100.5) rectangle ++(3.0,3);
\draw[wire] (94,90)--(r2hL1.south);
\draw[wire] (r2h02.north)--(111,52);
\draw[wire] (111,62)--(111,67);
\node[] (r2ellipsis2) at (111,71) {$\vdots$};
\draw[wire] (111,75)--(111,80);
\node[inner sep=0pt,minimum width=12mm,minimum height=3mm] (r2hL2) at (111,102) {};
\filldraw[draw=NNOutput,fill=NNOutput!18,line width=.5pt] (105.0,100.5) rectangle ++(3.0,3);
\filldraw[draw=NNOutput,fill=NNOutput!18,line width=.5pt] (108.0,100.5) rectangle ++(3.0,3);
\filldraw[draw=NNOutput,fill=NNOutput!18,line width=.5pt] (111.0,100.5) rectangle ++(3.0,3);
\filldraw[draw=NNOutput,fill=NNOutput!18,line width=.5pt] (114.0,100.5) rectangle ++(3.0,3);
\draw[wire] (111,90)--(r2hL2.south);
\node[stage,draw=NNOutput,fill=NNOutput!18,minimum width=22mm,minimum height=8mm] (r2head) at (77,117) {MLM头};
\draw[wire] (r2hL0.north)--(r2head.south);
\node[] (r2target) at (77,130) {猫};
\draw[wire] (r2head.north)--(r2target.south);
\node[nn heading] (r2sample) at (94,-12) {遮蔽样本2};
\node[nn heading] (rtitle) at (145,139) {共用层结构};
\node[] (rdin) at (145,17) {$\bm H^{(l-1)}$};
\node[stage,minimum width=34mm] (rd0) at (145,34) {双向多头注意力};
\draw[wire] (rdin.north)--(rd0.south);
\node[sum] (rd1) at (145,49) {$+$};
\draw[wire] (rd0.north)--(rd1.south);
\node[norm,minimum width=34mm] (rd2) at (145,62) {LN};
\draw[wire] (rd1.north)--(rd2.south);
\node[stage,minimum width=34mm] (rd3) at (145,82) {逐位置FFN};
\draw[wire] (rd2.north)--(rd3.south);
\node[sum] (rd4) at (145,97) {$+$};
\draw[wire] (rd3.north)--(rd4.south);
\node[norm,minimum width=34mm] (rd5) at (145,110) {LN};
\draw[wire] (rd4.north)--(rd5.south);
\node[] (rdout) at (145,124) {$\bm H^{(l)}$};
\draw[wire] (rd5.north)--(rdout.south);
\draw[wire] (145,24)--(167.0,24)--(167.0,49)--(rd1.east);
\fill[NNWire] (145,24) circle (.55);
\draw[wire] (145,70)--(167.0,70)--(167.0,97)--(rd4.east);
\fill[NNWire] (145,70) circle (.55);
\end{tikzpicture}
